\section{Rule-Based Models}
\begin{definition}[Rule based classifier]
A \textit{rule based classifier} is a model that uses a \textbf{rule set} of \texttt{if-then} rules to classify records.
\begin{equation}\label{ruleStructureClassifier}
    \overbrace{r_i}^{\textbf{Rule}} : \underbrace{(Cond_i)}_{\substack{\textbf{Antecedent}\\\text{or} \\\textbf{precondition}}} \to \overbrace{y_i}^{\textbf{Consequent}}\qquad \text{rule of a rule based classifier}
\end{equation}        
Each precondition is defined by a set of $k$ attribute-value pairs, or \textbf{conjunct}, such that:
    \begin{equation}
    Cond_i = \overbrace{(\underbrace{A_1 \text{ op } v_1}_{\textbf{conjunct}}) \land (A_2 \text{ op } v_2) \land \dots \land (A_k \text{ op } v_k)}^{\textbf{Conjuncts}}
    \end{equation}
$r_i$ \textbf{covers} an instance $x$ if the attributes of the instance satisfy the rule's antecedent.
\end{definition}
\begin{table}[h]
    \centering
    \begin{tabular}{|c|c|c|c|}
    \hline
        Name & Can Fly & Gives Birth & Blood Type \\
    \hline
    \hline
        Bat & Y & Y & W \\
    \hline
        Owl & Y & N & W \\
    \hline
        Platypus & N & N & W \\
    \hline
    \end{tabular}
    \caption{Small example dataset.}
    \label{tab:small_dataset}
\end{table}
\begin{example}
A rule based classifier classifies a test instance based on the rule triggered by the instance.
\begin{itemize}
    \item     $r_1 : (Can Fly = Y) \xrightarrow{} Bird$
    \item    $r_2 : (Gives Birth = Y) \land (Blood Type = W) \xrightarrow{} Mammal$
    \item     $r_3 : (Blood Type = C) \xrightarrow{} Reptile$
\end{itemize}
$Owl$ triggers $r_1$, so it is classified as a $Bird$. $Bat$ triggers both $r_1$ and $r_2$, a conflict. No rule covers $Platypus$, so there is no immediate way to assign it a class.
\end{example}

\begin{definition}
\textbf{Coverage of a rule:} fraction of records in the whole dataset that are covered by it, i.e., that satisfy the antecedent of a rule;
\begin{equation}
Coverage(r) = \frac{\overbrace{|A|}^{\text{number of records covered}}}{\underbrace{|D|}_{\text{total number of records}}} 
\end{equation}
\textbf{Accuracy of a rule} (sometimes called \textbf{precision})\textbf{:} among the records that satisfy the antecedent, the fraction that also satisfy the consequent.
\begin{equation}
    Accuracy(r) = \frac{\overbrace{|A \cap y|}^{\substack{\text{number of record}\\\text{ covered and} \\\text{correctly predicted}}}}{|A|}
\end{equation}
\end{definition}
\subsection{Properties of a Rule Set}

A rule set $R$ generated by a model can be evaluated using two key properties:
\begin{itemize}
    \item \textbf{Exhaustiveness:} Every combination of attribute values triggers at least one rule.
    \item \textbf{Mutual exclusivity:} Each record is covered by at most one rule in $R$.
\end{itemize}

Most rule-based classifiers lack one or both properties:
\begin{itemize}
    \item If $R$ is not exhaustive, a \textit{default rule} with an empty antecedent can classify uncovered instances.
    \item If $R$ is not mutually exclusive, the rules can be organized into an \textbf{ordered rule set} (or \textbf{decision list}), prioritizing rules by rank.
\end{itemize}

Rule ranking can follow:
\begin{itemize}
    \item \textbf{Rule-based ordering:} based on rule quality (e.g., accuracy).
    \item \textbf{Class-based ordering:} grouping rules with the same consequent.
\end{itemize}

For \textbf{unordered rule sets}, a \textbf{voting scheme} is used: each triggered rule casts a vote for its class. The predicted class is the one with the most votes, optionally weighted by rule confidence.

Unordered sets are less error-prone and cheaper to construct (no sorting), but more expensive at classification time, as each instance must be compared to all rules.

\subsection{Building a Rule Set}

Rule sets can be built via:
\begin{itemize}
    \item \textbf{Direct methods:} Rules are extracted directly from the data (e.g., RIPPER).
    \item \textbf{Indirect methods:} Rules are derived from other models (e.g., decision trees), typically ensuring mutual exclusivity and exhaustiveness.
\end{itemize}

\subsubsection{Direct Methods for Rule Extraction}

We focus on \textbf{RIPPER} (Repeated Incremental Pruning to Produce Error Reduction), a scalable rule learner that performs well on imbalanced and noisy datasets by using a validation set to prevent overfitting.

RIPPER employs a \textbf{sequential covering} strategy to build rules one class at a time:
\begin{itemize}
    \item For binary classification, the majority class becomes the default; rules are learned for the minority class.
    \item For multiclass tasks, classes are ordered by increasing frequency. Rules are built from the least frequent class to the most, with the last class as the default.
\end{itemize}

\begin{algorithm}
\caption{Sequential Covering Algorithm}
\begin{algorithmic}[1]
    \State $E \gets$ Training instances
    \State $A \gets$ Set of attribute-value pairs
    \State $Y_{\sigma} \gets \{y_1, y_2, \dots, y_k\}$ \Comment{Target classes}
    \State $R \gets \{\}$ \Comment{Initialize empty rule set}
    \For{$y \in Y_{\sigma} \setminus \{y_k\}$}
        \While{not stopping condition}
            \State $r \gets$ \Call{Learn-Rule}{$E, A, y$}
            \State $E \gets E \setminus \{x \in E : r \text{ covers } x\}$
            \State $R \gets R \cup \{r\}$
        \EndWhile
    \EndFor
    \State $R \gets R \cup \{\emptyset \rightarrow y_k\}$ \Comment{Default rule}
\end{algorithmic}
\end{algorithm}

\paragraph{Rule Learning Strategy}
RIPPER grows rules using a \textit{general-to-specific} approach, starting with an empty rule $r: \{\} \rightarrow +$ and iteratively adding conjuncts to maximize performance. The best conjunct is chosen using \textbf{FOIL's information gain}:
\begin{equation}
    FOIL_{gain} = p_1 \times \left( \log_2 \frac{p_1}{p_1 + n_1} - \log_2 \frac{p_0}{p_0 + n_0} \right)
\end{equation}
where $(p_0, n_0)$ and $(p_1, n_1)$ are the numbers of positive and negative examples before and after adding the conjunct.

\paragraph{Rule Pruning}
After growing a rule, RIPPER prunes it to enhance generalization, using:
\begin{equation}
    v = \frac{p - n}{p + n}
\end{equation}
where $p$ and $n$ are the positive and negative instances covered. If $v$ increases after pruning a conjunct, it is retained.

\paragraph{Model Selection with MDL}
To balance complexity and accuracy, RIPPER uses the \textbf{Minimum Description Length (MDL)} principle:
\begin{equation}
    Cost(M, D) = Cost(D|M) + \alpha \times Cost(M)
\end{equation}
where $M$ is the model, $D$ is the data, and $\alpha$ (default 0.5) controls the trade-off. Rule generation stops when the cost increases by at least $d$ bits (default $d = 64$).

\begin{definition}
\textbf{Pessimistic Error Estimate}:
\begin{equation}
    err_{gen}(T) = err(T) + \Omega \times \frac{k}{N_{\text{train}}}
\end{equation}
where $err(T)$ is the training error, $k$ is the number of rules, $N_{\text{train}}$ is the training set size, and $\Omega$ is the rule cost factor.
\end{definition}

\paragraph{Rule Set Construction}
RIPPER iteratively selects the best rule for the current positive class and removes all covered examples. Rule addition continues until the description length increases beyond the threshold.

\paragraph{Rule Set Optimization}
After initial rule generation, RIPPER optimizes the rule set. For each rule $r$, two alternatives are considered:
\begin{itemize}
    \item \textbf{Replacement rule} $r^*$: a new rule learned from scratch.
    \item \textbf{Revised rule} $r'$: an extended version of $r$ with additional conjuncts.
\end{itemize}
RIPPER selects the alternative (or retains $r$) that minimizes the MDL cost.
\subsubsection{Indirect Methods for Rule Extraction}

Indirect methods derive rule sets from other models, typically unpruned decision trees. Each path from the root to a leaf becomes a classification rule: the conjunction of conditions along the path forms the antecedent, and the majority class at the leaf forms the consequent.

A well-known approach is \textbf{C4.5rules}, which creates one rule per path. For each rule $r: A \rightarrow y$, simplified versions $r': A' \rightarrow y$ are generated by removing conjuncts from $A$. The version with the lowest pessimistic error is retained if it improves upon the original. Duplicate rules are eliminated, and simplification continues until no further improvements are possible.

After generating rules, C4.5rules applies \textbf{class-based ordering}, grouping rules by predicted class. Rule subsets are ranked by their description length, with the class having the shortest total description length given highest priority. This assumes the best rules are found in the most concise subset.\footnote{While RIPPER prioritizes precision and learns from the least frequent classes first, C4.5rules evaluates subsets independently using pessimistic error and selects based on estimated quality.}

\subsection{Characteristics of Rule-Based Models}

Rule-based classifiers, like decision trees, partition the input space using rectilinear decision boundaries. Unlike trees, where each instance follows a single path, rule-based models may trigger multiple rules for the same instance, enabling them to approximate more complex functions.

Both approaches support categorical and continuous features and can handle binary or multiclass tasks. Rule-based classifiers offer:
\begin{itemize}
    \item High interpretability
    \item Competitiveness with decision trees in performance
    \item Robustness to redundant features (often selecting only one among correlated attributes)
    \item Natural handling of class imbalance (e.g., RIPPER prioritizes minority classes)
\end{itemize}

They also have two clear limitations:
\begin{itemize}
    \item They cannot handle missing values in test data.
    \item They can produce conflicting predictions when several rules apply at once.
\end{itemize}

\noindent
\textbf{Summary:} Decision trees are typically preferred due to their robustness to missing data. Rule-based classifiers are advantageous in highly imbalanced settings but struggle with missing values during inference.
